\chapter*{General Introduction}
\addcontentsline{toc}{chapter}{General Introduction}
\markboth{General Introduction}{General Introduction}
\label{chap:general-introduction}

The transition toward a low-carbon economy has become a global priority. While
renewable electricity generation has expanded rapidly, several sectors, heavy
industry, aviation, maritime transport, and high-temperature industrial
processes, cannot be decarbonized through direct electrification alone. For
these hard-to-abate sectors, the \textbf{Power-to-X} (PtX) paradigm has emerged
as a promising pathway: renewable electricity is converted into storable
carriers and molecules, such as green hydrogen and its derivatives (ammonia,
methanol, and sustainable aviation fuel), that can be transported and used where
electrons cannot.

Power-to-X projects are capital-intensive and structurally complex. Their
viability depends on the interplay of renewable generation, electrolysis,
storage, market prices, off-take agreements, and financing assumptions over a
project's lifetime. At the \emph{pre-feasibility} stage, sponsors must compare
technical, operational, and economic trade-offs while many assumptions remain
uncertain, which calls for robust and reproducible analytical tools. To support
this process, BCG~X, the technology, design, and engineering unit of Boston
Consulting Group (BCG), developed the \textbf{Energy Transition Optimizer}, an
internal product that simulates the operation of a candidate PtX project and
produces auditable financial indicators. Built in-house and then adapted for
several clients, it rests on a deterministic optimization engine: a project is
described as a typed graph of assets, conservation nodes, and markets, together
with hundreds of parameters, from which the engine derives hour-by-hour
operating schedules and economic indicators such as net present value, internal
rate of return, and the levelized cost of the output product.

As the product matured, two obstacles limited its value. First, it sat on a
\textbf{fragile platform}, with an unreliable delivery pipeline, infrastructure
drift, security findings in a self-hosted identity layer, a dated interface, and
weak operational observability. Second, and more fundamentally, it was
\textbf{expert-only}: obtaining an answer required hand-building a valid
topology, binding assumptions to every node, configuring parameter sweeps, and
reading a raw population of results. At the same time, advances in large
language models opened the possibility of making such tools accessible through
natural language. Integrating that capability into a product whose value rests
on the reproducibility and auditability of its computations, however, raises a
central tension: automation must \emph{assist} the analyst without becoming the
source of numerical truth.

It is in this context that I carried out my end-of-study internship at BCG~X,
within the Energy Transition Optimizers product team. Conducted as a team
effort, the work pursued two complementary goals: consolidating the platform
foundation, and introducing a \textbf{governed multi-agent conversational
layer}. In this layer, a director agent routes each request to specialized
agents (data ingestion, topology design, scenario exploration, and result
interpretation), while every state-changing action remains a typed proposal that
is validated and then subject to explicit human approval, so that the
deterministic engine remains the sole authority over numerical results. Three
components complete this layer and receive particular attention in this report: the
\textbf{result-interpretation agent}, a read-only sub-agent that explains the
economics of a simulation population in natural language while keeping every
figure traceable to a deterministic result; the associated
\textbf{visualization layer}; and an \textbf{evaluation harness} that measures
the factual grounding and financial soundness of the explanations produced.


This report is organized into five chapters. Chapter~\ref{chap:context} presents
the host organization, the energy-transition context, the product, and the
problem the internship addresses. Chapter~\ref{chap:background} reviews the
technical background and the state
of the art, from techno-economic optimization to large-language-model agents and
their evaluation, and justifies the chosen direction.
Chapter~\ref{chap:requirements} analyzes and
specifies the functional and non-functional requirements.
Chapter~\ref{chap:design} details the
design and architecture of the modernized platform and its governed multi-agent
layer, with particular attention to the interpretation, visualization, and
evaluation components. Chapter~\ref{chap:implementation} describes the
implementation and reports the
evaluation results and measured impact. A short unnumbered chapter then records
the United States patent granted on the optimizer's core method while this work
was under way, before a general conclusion summarizes the
outcomes and outlines future perspectives.

